\begin{resumo}[Resumo Estendido]
\begin{otherlanguage*}{brazil}

\section*{1. Contextualização e Motivação}

Com a expansão contínua da sociedade digital e o avanço da computação em nuvem, aprendizado de máquina e internet das coisas, a geração global de dados cresce a taxas que desafiam a infraestrutura de armazenamento convencional. Mídias magnéticas e ópticas tradicionais apresentam limitações estruturais intrínsecas: vida útil operacional reduzida (geralmente entre 5 e 10 anos), exigência constante de ciclos de migração de dados e elevado consumo contínuo de energia elétrica para manutenção de campos magnéticos e refrigeração de centros de dados.

Nesse cenário, o ácido desoxirribonucleico (DNA) sintético desponta como uma das mídias mais promissoras para arquivamento digital de longo prazo (\emph{cold data storage}). O DNA oferece uma densidade de informação teórica de até 455 exabytes por grama, durabilidade passiva de milênios quando protegido contra umidade e temperaturas extremas, além de estabilidade físico-química sem necessidade de alimentação energética contínua. Contudo, o canal molecular de armazenamento em DNA impõe restrições biofísicas e estocásticas severas: a síntese química e o sequenciamento em larga escala são propensos a erros de substituição, inserção, deleção e, criticamente, ao apagamento completo de oligonucleotídeos. Além disso, sequências com teores extremos de Guanina-Citosina (GC) ou com repetições homopoliméricas longas inviabilizam a amplificação por PCR e o sequenciamento homogêneo.

Para superar esses obstáculos, \citeonline{erlich2017} propuseram a arquitetura DNA Fountain, aliando códigos de fonte inesgotável (\emph{fountain codes} / códigos LT) a filtros de restrições biológicas (\emph{screening}). Inspirado por esse avanço e motivado pela carência de ferramentas didáticas, modulares e estritamente reproduzíveis no contexto acadêmico de Engenharia de Software no Brasil, este Trabalho de Conclusão de Curso concebeu, implementou e avaliou um simulador modular do pipeline DNA Fountain em software puro.

\section*{2. Formulação do Problema e Objetivos}

A questão central de pesquisa investigada neste trabalho é: \emph{como projetar e validar um simulador modular do pipeline DNA Fountain que represente a codificação LT, restrições biológicas, perdas por apagamento e a decodificação por peeling, permitindo a análise experimental sistemática do impacto dos parâmetros $c$, $\delta$, número de blocos-fonte $K$, probabilidade de perda e overhead na taxa de recuperação e no custo computacional?}

O objetivo geral do trabalho foi desenvolver e avaliar um simulador modular do pipeline DNA Fountain, implementando todas as suas etapas lógicas essenciais em Python puro (biblioteca padrão). Os objetivos específicos englobaram: (1)~formalizar matematicamente os códigos LT, a Distribuição Robusta de Soliton (RSD) e o decodificador \emph{peeling}; (2)~modelar computacionalmente o fluxo de segmentação, síntese de gotas com cabeçalhos baseados em sementes e conversão nucleotídica; (3)~implementar o simulador com separação estrita de responsabilidades; (4)~incorporar filtros de restrição biológica para conteúdo GC ($40\%$ a $60\%$) e limite de homopolímeros ($\le 3$ bases); (5)~simular um canal de apagamento estocástico independente (abstração BEC); (6)~instrumentar a coleta de métricas de sucesso, dinâmica de fila do \emph{ripple}, operações XOR e tempos de execução; (7)~avaliar experimentalmente uma grade fatorial ampla de configurações em regime de blocos finitos ($K \le 64$); e (8)~garantir estrita reprodutibilidade determinística por meio do isolamento de domínios pseudo-aleatórios.

\section*{3. Fundamentação Teórica: Códigos LT e DNA Fountain}

Diferentemente dos códigos de correção de erros com taxa fixa clássicos (como Reed--Solomon), os códigos \emph{fountain} operam sem taxa de codificação estática pré-determinada. O codificador particiona o arquivo em $K$ blocos-fonte de tamanho fixo $L$ bytes e é capaz de sintetizar um fluxo potencialmente ilimitado de pacotes codificados (\emph{gotas} ou \emph{droplets}). Cada gota é gerada por meio de três passos: (a) amostragem de um grau inteiro $d \in \{1, \ldots, K\}$ a partir de uma distribuição de probabilidades; (b) seleção uniforme e sem reposição de $d$ blocos-fonte vizinhos; e (c) combinação dos blocos selecionados através de somas bit a bit no corpo finito binário $GF(2)$ via operações XOR ($\oplus$).

A distribuição de graus é o fator determinante do desempenho do código. \citeonline{luby2002} introduziu a Distribuição Ideal de Soliton $\rho(d)$ e, para sanar sua vulnerabilidade à interrupção prematura em instâncias finitas, a Distribuição Robusta de Soliton (RSD) $\mu(d) = (\rho(d) + \tau(d)) / \beta$, que adiciona o termo corretivo $\tau(d)$ governado pelos parâmetros $c > 0$ e $\delta \in (0, 1)$, pela escala $R = c \ln(K/\delta)\sqrt{K}$ e pelo índice modal $m = \lfloor K/R \rfloor$. Em regimes finitos com blocos curtos ($K \le 64$) e valores elevados de $c$, o simulador incorpora uma salvaguarda para o caso de borda em que $m_{\text{raw}} = \lfloor K/R \rfloor < 1$, fixando $m = 1$ e alocando toda a massa corretiva em grau 1 para assegurar a consistência estocástica do gerador.

A decodificação é realizada pelo algoritmo iterativo de \emph{peeling}, que propaga crenças em tempo linear sobre o grafo bipartido conectando blocos-fonte e gotas recebidas. Quando uma gota possui grau residual 1, seu conteúdo revela diretamente o bloco-fonte correspondente. Esse bloco é então removido (via XOR) de todas as gotas vizinhas incidentes no grafo, reduzindo seus graus residuais. Gotas cujo grau residual atinge 1 entram em uma fila operacional (\emph{ripple}). O processo de decodificação é bem-sucedido quando todos os $K$ blocos são reconstruídos antes que a fila se esvazie. O esgotamento prematuro da fila antes da resolução completa caracteriza a interrupção por \emph{ripple die-out}.

\section*{4. Arquitetura de Software do Simulador}

O simulador foi concebido sob princípios rigorosos de Engenharia de Software, priorizando alta coesão, baixo acoplamento e dependência exclusiva da biblioteca padrão do Python 3.12 (CPython). A arquitetura é encapsulada no pacote \texttt{dna\_fountain}, estruturado em módulos especializados:
\begin{itemize}
    \item \texttt{config.py}: centraliza os hiperparâmetros do sistema por meio de \emph{dataclasses} tipadas (\texttt{DNAFountainConfig}, \texttt{ExperimentConfig});
    \item \texttt{preprocessing.py}: executa o particionamento de arquivos binários em $K$ blocos de tamanho $L$ com preenchimento determinístico de zeros (\emph{zero-padding});
    \item \texttt{soliton.py}: calcula as probabilidades da RSD, normaliza a distribuição cumulativa (CDF) e realiza a amostragem determinística de graus e vizinhanças a partir de sementes inteiras;
    \item \texttt{screening.py}: implementa a validação físico-química, inspecionando o teor GC e o comprimento máximo de homopolímeros nas sequências candidatas;
    \item \texttt{lt\_codec.py}: orquestra a codificação (\texttt{DNAFountainEncoder}) e a decodificação \emph{peeling} (\texttt{DNAFountainDecoder}) com indexação invertida $O(1)$ em tabelas hash e instrumentação de passos;
    \item \texttt{dna\_conversion.py}: realiza a conversão bijetiva entre bytes e nucleotídeos à taxa bruta de 2 bits por base ($A \leftrightarrow 00$, $C \leftrightarrow 01$, $G \leftrightarrow 10$, $T \leftrightarrow 11$);
    \item \texttt{channel.py}: modela o canal de apagamento de pacotes independentes com probabilidade $p_{\text{loss}}$;
    \item \texttt{experiments.py}: automatiza a execução de lotes paramétricos e consolida métricas em arquivos CSV estruturados.
\end{itemize}

O enquadramento binário de cada oligonucleotídeo sintético compreende 7 bytes (28 bases): 2 bytes para o identificador da semente (máscara XOR \texttt{0x1B1B}), $L = 4$ bytes de carga útil e 1 byte de checksum CRC-8. O uso da semente permite ao decodificador regenerar exatamente os mesmos índices de vizinhança sem transmitir metadados explícitos de conectividade do grafo.

A reprodutibilidade estrita é assegurada pelo isolamento de três fluxos pseudo-aleatórios independentes: o PRNG da semente de 16 bits para conectividade da gota, o PRNG estocástico do canal de apagamento e o PRNG mestre de orquestração dos lotes de teste. A suite de testes automatizados (\texttt{pytest}) conta com 10 procedimentos de validação canônica cobrindo todas as primitivas do sistema.

\section*{5. Metodologia Experimental}

O protocolo experimental investigou uma grade fatorial de 216 combinações paramétricas formadas pelo cruzamento de: número de blocos-fonte $K \in \{16, 32, 64\}$; parâmetro $c \in \{0.1, 0.3, 0.5, 0.8\}$; parâmetro $\delta \in \{0.01, 0.10\}$; overhead nominal $\epsilon \in \{0.25, 0.50, 1.00\}$; e probabilidade de perda no canal $p_{\text{loss}} \in \{0.0, 0.10, 0.20\}$. Cada ponto fatorial foi submetido a 100 repetições independentes, totalizando $21.600$ simulações completas. Um experimento adicional independente de amostragem avaliou $2.000$ oligonucleotídeos para isolar a seletividade físico-química do módulo de screening em $K=32$.

As métricas coletadas incluíram: taxa de sucesso de decodificação ($P_{\text{succ}}$) com intervalos exatos de confiança de Clopper--Pearson a 95\%; contagem total de operações XOR durante a decodificação; tamanho máximo ($Q_{\max}$) e mínimo ($Q_{\min}$) da fila de \emph{ripple}; taxa de rejeição no screening ($R_{\text{screen}}$); e tempos de processamento computacional em milissegundos. Testes exatos bicaudais de Fisher e testes de qui-quadrado de aderência foram aplicados para determinar a significância estatística das hipóteses avaliadas.

\section*{6. Resultados e Discussão}

Os principais achados experimentais organizam-se em seis dimensões:
\begin{enumerate}
    \item \textbf{Validação Funcional de Referência}: sob canal ideal ($p_{\text{loss}} = 0.0$), $\delta = 0.10$ e overhead $\epsilon = 1.00$, as configurações $c = 0.1$ sustentaram $100\%$ de sucesso para $K = 32$ e $K = 64$ (IC 95\%: $[96.4\%, 100.0\%]$). Para $K = 16$, observou-se taxa de sucesso de $90\%$; a auditoria detalhada das 10 falhas revelou que 9 decorreram de interrupção prematura da fila de peeling em conjuntos de parada (\emph{stopping sets}) onde todas as gotas restantes apresentavam grau residual $\ge 2$, e apenas 1 decorreu de ausência inicial de cobertura de bloco no grafo.
    
    \item \textbf{Impacto do Screening Biológico}: o filtro físico-químico rejeitou $40.2\%$ dos candidatos aleatórios gerados. O teste qui-quadrado de aderência revelou que o grau médio das gotas aceitas ($4.765$) não diferiu significativamente da expectativa teórica da RSD ($4.945$, $\chi^2 = 3.64$, $p \approx 0.30$), demonstrando que o descarte de sequências inadequadas não distorce a conectividade média do grafo.
    
    \item \textbf{Dinâmica do Ripple}: a evolução da fila confirmou que regimes com recuperação bem-sucedida atingem picos médios de fila adequados ($Q_{\max} \approx 10.5$ para $K=32$, $\epsilon=0.50$ com $c=0.1$). Em contrapartida, configurações patológicas com $c \ge 0.5$ concentraram excessivamente a massa em grau 1 ($m=1$, representando $\approx 73\%$ dos pacotes), privando o grafo de arestas intermediárias e resultando em $0\%$ de sucesso (0/100).
    
    \item \textbf{Transição de Overhead}: as curvas de sucesso em função do overhead tornam-se sensivelmente mais íngremes à medida que $K$ aumenta de 16 para 64, em plena concordância com a teoria assintótica dos códigos LT.
    
    \item \textbf{Efeito dos Parâmetros $c$ e $\delta$}: para $K = 16$ com orçamento de redundância restrito ($\epsilon = 0.25$), a configuração $c = 0.3$ superou estatisticamente $c = 0.1$ ($36.0\%$ vs. $8.0\%$, teste exato de Fisher, $p = 2.2 \times 10^{-6}$). Esse benefício decorre do deslocamento de $m$ de 7 para 2, concentrando expressivos $70.1\%$ da probabilidade em gotas de grau 2 ($\mu(2) \approx 0.701$), que fornecem equações binárias prontas para alimentar a fila assim que o primeiro bloco unitário é resolvido. Para $K=32$, a diferença amostral entre $c=0.3$ e $c=0.1$ ($38.0\%$ vs. $30.0\%$) não apresentou significância estatística ($p \approx 0.30$).
    
    \item \textbf{Resiliência a Apagamentos no Canal}: sob perdas estocásticas de $10\%$ e $20\%$, o sistema preservou $100\%$ de recuperação para $K=32$ e $K=64$ desde que o overhead nominal compensasse a fração de perdas ($\epsilon = 1.00$), comprovando a flexibilidade da abordagem \emph{fountain} sem necessidade de protocolos de retransmissão por canal de retorno.
\end{enumerate}

\section*{7. Conclusões e Trabalhos Futuros}

Este trabalho concluiu com êxito a modelagem, implementação e caracterização experimental do pipeline DNA Fountain sob a perspectiva da Engenharia de Software. Foi demonstrado que, no regime de blocos finitos, o sucesso da decodificação por peeling depende criticamente da calibração entre a distribuição de graus, a contagem de blocos e a margem de redundância, com destaque para a relevância das gotas de grau 2 em orçamentos restritos.

O código-fonte completo, os scripts de execução em lote e as matrizes de dados estão disponíveis sob licença de código aberto em repositório público (\url{https://github.com/GabrielSMonteiro/DNA_Fountain}). Como direções futuras de pesquisa, destacam-se a modelagem de canais com ruído de fase (inserções e deleções via alinhamento de Levenshtein ou códigos HEDGES), a incorporação de códigos corretores concatenados (Reed--Solomon interno de 2 bytes), a avaliação de distribuições alternativas para códigos Raptor e BATS, e a prototipagem de aceleração do decodificador peeling em hardware dedicado (FPGA Digilent Basys 3).

\vspace{0.5cm}
\textbf{Palavras-chave}: Armazenamento de dados em DNA. DNA Fountain. Códigos LT. Distribuição Robusta de Soliton. Decodificador Peeling. Simulação de Software.

\end{otherlanguage*}
\end{resumo}
